
\usepackage{polyglossia}

\usepackage{fontspec}

\usepackage{ifplatform}

\ifwindows
  \newfontfamily{\cyrillicfont}{Times New Roman}
  \setmainfont{Times New Roman}
  \newfontfamily{\cyrillicfonttt}{Courier New}
  \setmonofont{Courier New}
  \newfontfamily{\cyrillicfontsf}{Microsoft Sans Serif}[Script=Cyrillic]
  \setsansfont{Microsoft Sans Serif}
\else
  \newfontfamily\russianfont{Linux Libertine O}[Script=Cyrillic]
  \newfontfamily\cyrillicfontsf{Linux Biolinum O}[Script=Cyrillic]
  \newfontfamily\cyrillicfonttt{DejaVu Sans Mono}[Script=Cyrillic]
  \setmainfont{Linux Libertine O}
  \setsansfont{Linux Biolinum O}
  \setmonofont{DejaVu Sans Mono}[Script=Cyrillic,Scale=MatchLowercase]
\fi

\setmainlanguage[numerals=cyrillic]{russian}
\setotherlanguages{english}
\usepackage{csquotes}

\usepackage{xunicode} % some extra unicode support
%\usepackage[utf8x]{inputenc}
\usepackage{xltxtra} % \XeLaTeX macro
\defaultfontfeatures{Ligatures=TeX}

%\setromanfont{Charis SIL}
%\setsansfont{Liberation Sans}
%\setmonofont{PT Mono}
%\setmainfont{Liberation Serif} % this allows to use sans-serif as default font

%нумерация справа и колонтитулы справа вверху
\usepackage{fancyhdr}
\usepackage[a4paper,left=30mm,right=15mm,top=20mm,bottom=20mm,bindingoffset=0cm]{geometry}%

\usepackage{indentfirst}


\usepackage[protrusion=true]{microtype}
\emergencystretch 3em

\usepackage{mathtools}
\usepackage{amsfonts}
\usepackage{amssymb}
\usepackage{amsmath}
\usepackage{amsthm}

\usepackage{calc}
\usepackage{ifthen}
\usepackage{graphicx}
\usepackage{tabularx}
\usepackage{array}
\usepackage{ragged2e}
\newcolumntype{P}[1]{>{\RaggedRight\arraybackslash}p{#1}}
\usepackage{pdfpages}
\usepackage{longtable}
\usepackage{multirow}
\usepackage{float}
\usepackage{placeins}
\makeatletter %попытка избавиться от нежелательных разрывов строк, вроде не очень работает
\def\nobreakhline{%
  \noalign{\ifnum0=`}\fi
    \penalty\@M
    \futurelet\@let@token\LT@@nobreakhline}
\def\LT@@nobreakhline{%
  \ifx\@let@token\hline
    \global\let\@gtempa\@gobble
    \gdef\LT@sep{\penalty\@M\vskip\doublerulesep}% <-- change here
  \else
    \global\let\@gtempa\@empty
    \gdef\LT@sep{\penalty\@M\vskip-\arrayrulewidth}% <-- change here
  \fi
  \ifnum0=`{\fi}%
  \multispan\LT@cols
     \unskip\leaders\hrule\@height\arrayrulewidth\hfill\cr
  \noalign{\LT@sep}%
  \multispan\LT@cols
     \unskip\leaders\hrule\@height\arrayrulewidth\hfill\cr
  \noalign{\penalty\@M}%
  \@gtempa}
\makeatother

\usepackage[unicode=true]{hyperref}
\usepackage{xurl}
\usepackage{color}
\usepackage{pgf}
\usepackage{pstheorems}

\usepackage{titling}

% Настройка списков (без лишних вертикальных отступов)
\usepackage{paralist}
\setdefaultenum{1.}{1.}{1.}{1.}
\setdefaultitem{--}{}{}{}
%\setlength\itemsep{-1em}
\let\itemize\compactitem
\let\enditemize\endcompactitem
\let\enumerate\compactenum
\let\endenumerate\endcompactenum
\let\description\compactdesc
\let\enddescription\endcompactdesc
\pltopsep=\smallskipamount
\plitemsep=0pt
\plparsep=0pt
% Команда для отмены разрыва страниц перед списками
\makeatletter 
\newcommand\mynobreakpar{\par\nobreak\@afterheading} 
\makeatother
%%%%%%

\usepackage[singlelinecheck=true,labelsep=endash]{caption}
\usepackage{subcaption}
\captionsetup[table]{justification=justified}
\captionsetup[figure]{justification=justified,name=Рисунок,singlelinecheck=on,font=onehalfspacing}

\usepackage{titlesec}
%\titleformat{\sectionname}[runin]{начертание}{счетчик, например \thesubparagraph}{0pt}{}{}

\titleformat{\chapter}[block]{\centering\normalfont\Large\bfseries}{\thechapter.}{1ex}{}{}
\titlespacing{\chapter}{0pt}{0em}{2em}

\titleformat{\section}[block]{\normalfont\large\bfseries}{\thesection}{1ex}{}{}
\titlespacing{\section}{0pt}{0em}{1ex}

\titleformat{\subsection}[block]{\normalfont\normalsize\bfseries}{\thesubsection}{1ex}{}{}
\titlespacing{\section}{0pt}{0em}{1ex}

	% paragraph и subparagraph -- в тексте, без отступов
% \titlespacing{command}{left spacing}{before spacing}{after spacing}[right]

\titleformat{\paragraph}[runin]{\normalfont\normalsize\bfseries}{\theparagraph}{0pt}{}{}
\titlespacing{\paragraph}{0pt}{0em}{0ex} %не добавлять дополнительных отступов
% \titlespacing{\paragraph}{\parindent}{0em}{10ex}[1em] %добавить отступы

\titleformat{\subparagraph}[runin]{\normalfont\normalsize\bfseries}{\thesubparagraph}{0pt}{}{}
\titlespacing{\subparagraph}{0pt}{0em}{0ex}


% Своё название для Cписка литературы
\usepackage[title, titletoc]{appendix}
\addto\captionsrussian{% Replace "english" with the language you use
	\renewcommand{\contentsname}%
	{Содержание}%
}

%\renewcommand{\appendixname}{Приложение}% Change "chapter name" for Appendix chapters
%\renewcommand{\cftchapdotsep}{\cftdotsep}

\usepackage{mathpartir}

\makeatletter
\let\ps@plain\ps@fancy              % Подчиняем первые страницы каждой главы общим правилам
\makeatother
\pagestyle{fancy}
\fancyhf{}
\fancyfoot[C]{\thepage}
\renewcommand{\headrulewidth}{0pt}
\renewcommand{\footrulewidth}{0pt}
\renewcommand{\baselinestretch}{1.5}
\newcommand{\headertext}[1]{\fancyhead[R]{\tiny{#1}}}


% поддержка нумерованных списков с нумерацией в виде русских букв
\usepackage{enumitem}
\makeatletter
\AddEnumerateCounter{\asbuk}{\russian@asbuk@alph}{а} % section 8.1
\renewcommand{\thesubfigure}{\asbuk{subfigure}}
\makeatother

%% Список литературы

\usepackage[
  style=gost-numeric,
  sorting=none,
  language=auto,
  autolang=other
]{biblatex}
\addbibresource{chapters/biblio.bib}
\renewcommand\multicitedelim{,\space}

\usepackage{tikz}
\usetikzlibrary{shapes.geometric,arrows.meta,positioning}
\usepackage{hhline}

%\frenchspacing %% изменение расстояние до и после точек в ряде случаев

\renewcommand{\theenumi}{\arabic{enumi}}
\renewcommand{\theenumii}{\arabic{enumii}}
\renewcommand{\theenumiii}{\arabic{enumiii}}
\renewcommand{\theenumiv}{\arabic{enumiv}}

\renewcommand{\labelenumi}{\theenumi.}
\renewcommand{\labelenumii}{\theenumi.\theenumii.}
\renewcommand{\labelenumiii}{\theenumi.\theenumii.\theenumiii.}
\renewcommand{\labelenumiv}{\theenumi.\theenumii.\theenumiii.\theenumiv.}

%\newenvironment{annotation}{\textbf{Аннотация.} \textit}{}
\theoremstyle{plain}
\newtheorem*{annotation}{Аннотация}

\makeatletter
\newcommand*{\projecttypefulldative}[1]{\gdef\@projecttypefulldative{#1}}
\newcommand*{\theprojecttypefulldative}{\@projecttypefulldative}
\newcommand*{\projecttypeshort}[1]{\gdef\@projecttypeshort{#1}}
\newcommand*{\theprojecttypeshort}{\@projecttypeshort}
\newcommand*{\authorfulldative}[1]{\gdef\@authorfulldative{#1}}
\newcommand*{\theauthorfulldative}{\@authorfulldative}
\newcommand*{\authorgroup}[1]{\gdef\@authorgroup{#1}}
\newcommand*{\theauthorgroup}{\@authorgroup}
\newcommand*{\supervisor}[1]{\gdef\@supervisor{#1}}
\newcommand*{\thesupervisor}{\@supervisor}
\newcommand*{\consultant}[1]{\gdef\@consultant{#1}}
\newcommand*{\theconsultant}{\@consultant}
\newcommand{\projecttasks}[1]{\gdef\@projecttasks{#1}}
\newcommand{\theprojecttasks}{\@projecttasks}
\newcommand{\projecttask}[5]{#1 & #2 & #3 & #4 & #5 \\\hline}
\newcommand*{\taskliterature}[1]{\gdef\@taskliterature{#1}}
\newcommand*{\thetaskliterature}{\@taskliterature}
\newcommand*{\taskdate}[1]{\gdef\@taskdate{#1}}
\newcommand*{\thetaskdate}{\@taskdate}
\newcommand*{\supervisortaskapproval}[1]{\gdef\@supervisortaskapproval{#1}}
\newcommand*{\thesupervisortaskapproval}{\@supervisortaskapproval}
\newcommand*{\authortaskapproval}[1]{\gdef\@authortaskapproval{#1}}
\newcommand*{\theauthortaskapproval}{\@authortaskapproval}
\newcommand*{\authorrspzapproval}[1]{\gdef\@authorrspzapproval{#1}}
\newcommand*{\theauthorrspzapproval}{\@authorrspzapproval}
\newcommand*{\supervisorrspzapproval}[1]{\gdef\@supervisorrspzapproval{#1}}
\newcommand*{\thesupervisorrspzapproval}{\@supervisorrspzapproval}
\newcommand*{\consultantrspzapproval}[1]{\gdef\@consultantrspzapproval{#1}}
\newcommand*{\theconsultantrspzapproval}{\@consultantrspzapproval}
\newcommand*{\supervisorrspzgrade}[1]{\gdef\@supervisorrspzgrade{#1}}
\newcommand*{\thesupervisorrspzgrade}{\@supervisorrspzgrade}
\newcommand*{\authorpzapproval}[1]{\gdef\@authorpzapproval{#1}}
\newcommand*{\theauthorpzapproval}{\@authorpzapproval}
\newcommand*{\supervisorpzapproval}[1]{\gdef\@supervisorpzapproval{#1}}
\newcommand*{\thesupervisorpzapproval}{\@supervisorpzapproval}
\newcommand*{\consultantpzapproval}[1]{\gdef\@consultantpzapproval{#1}}
\newcommand*{\theconsultantpzapproval}{\@consultantpzapproval}
\newcommand*{\supervisorpzgrade}[1]{\gdef\@supervisorpzgrade{#1}}
\newcommand*{\thesupervisorpzgrade}{\@supervisorpzgrade}
\makeatother

\newcommand{\sign}[3][0pt]{%
  \tikz[overlay]{\node[yshift=#1]{\includegraphics[scale=#2]{img/signatures/#3.png}}}%
}
\newcommand{\signat}[4][0pt]{%
  \begin{tikzpicture}[overlay]
    \node[xshift=-10pt,yshift=#1](c){\includegraphics[scale=#2]{img/signatures/#3.png}};
    \node[xshift=10pt,yshift=-10pt]{\scriptsize\textit{#4}};
  \end{tikzpicture}
}

\newcommand{\emptyfield}{\tikz[overlay]{\draw[thin,yshift=-1.28ex](0,0)--(5,0)}}

\newcounter{projecttasknumber}
\newcommand{\projecttasknum}{\setcounter{projectsubtasknumber}{0}\stepcounter{projecttasknumber}\theprojecttasknumber.}

\newcounter{projectsubtasknumber}
\newcommand{\projectsubtasknum}{\stepcounter{projectsubtasknumber}\theprojecttasknumber.\theprojectsubtasknumber.}

\usepackage{listings}

\renewcommand{\lstlistingname}{Листинг}

\lstset{
  basicstyle=\ttfamily\small,
  tabsize=2,
  showstringspaces=false,
  columns=fullflexible,
  keepspaces=true,
  numbers=left,
  stepnumber=1,
  numberstyle=\tiny\color{gray},
  breaklines=true,
  breakatwhitespace=true,
  framesep=6pt,
  aboveskip=0.5em,
  belowskip=0.5em,
  lineskip=0pt,
  abovecaptionskip=1em,
  captionpos=t,
  extendedchars=true,
  literate={Ö}{{\"O}}1
  {Ä}{{\"A}}1
  {Ü}{{\"U}}1
  {ß}{{\ss}}1
  {ü}{{\"u}}1
  {ä}{{\"a}}1
  {ö}{{\"o}}1
  {~}{{\textasciitilde}}1
  {…}{\ldots}1
  {–}{-}1
  {\ }{ }1
}

\headertext{}

\authorgroup{TODO: группа}
\author{TODO: Фамилия И. О.}
\authorfulldative{TODO: Фамилии Имени Отчеству}
\supervisor{TODO: руководитель}
\consultant{\emptyfield}
\projecttypefulldative{выпускной квалификационной работе}
\projecttypeshort{ВКР}

%%% Local Variables:
%%% TeX-engine: xetex
%%% eval: (setq-local TeX-master (concat "../" (seq-find (-cut string-match ".*-1-task\.tex$" <>) (directory-files ".."))))
%%% End:

\title{Разработка сервиса автоматизации отбора кандидатов\\
  на основе обработки естественного языка}

\taskdate{TODO: дата выдачи задания}

\projecttasks{%
  \projecttask{\bfseries\projecttasknum}{\bfseries Аналитическая часть}{}{}{}%
  \projecttask{\projectsubtasknum}%
  {Выполнить анализ предметной области рекрутмента, существующих подходов к автоматизации отбора кандидатов и ограничений ручной обработки резюме и вакансий.}%
  {Аналитический раздел ПЗ, список источников}%
  {TODO: срок}{}%
  \projecttask{\projectsubtasknum}%
  {Сформулировать требования к целевому сервису и сценарии взаимодействия пользователей с системой.}%
  {Функциональные требования, Use Case, бизнес-модели AS IS и TO BE}%
  {TODO: срок}{}%
  \projecttask{\bfseries\projecttasknum}{\bfseries Теоретическая часть}{}{}{}%
  \projecttask{\projectsubtasknum}%
  {Разработать модель сопоставления резюме и вакансий на основе методов обработки естественного языка и векторного представления текстов.}%
  {Описание модели, критерии ранжирования кандидатов}%
  {TODO: срок}{}%
  \projecttask{\projectsubtasknum}%
  {Подготовить методику экспериментальной проверки модели на тестовых данных.}%
  {Описание датасета, метрик и ожидаемых результатов}%
  {TODO: срок}{}%
  \projecttask{\bfseries\projecttasknum}{\bfseries Инженерная часть}{}{}{}%
  \projecttask{\projectsubtasknum}%
  {Спроектировать архитектуру сервиса, включая компоненты обработки резюме, обработки вакансий, сопоставления кандидатов и хранения данных.}%
  {Архитектурное описание, диаграммы, выбор технологий}%
  {TODO: срок}{}%
  \projecttask{\projectsubtasknum}%
  {Реализовать прототип сервиса и проверить работоспособность выбранного подхода на тестовых сценариях.}%
  {Исходный код, результаты проверки, описание ограничений прототипа}%
  {TODO: срок}{}%
  \projecttask{\bfseries\projecttasknum}{\bfseries Технологическая и практическая часть}{}{}{}%
  \projecttask{\projectsubtasknum}%
  {Оформить пояснительную записку и иллюстративный материал для защиты ВКР.}%
  {Текст ПЗ, презентация, приложения}%
  {TODO: срок}{}%
  \projecttask{\bfseries\projecttasknum}%
  {\bfseries Заменить временные титульные материалы официальными бланками кафедры перед сдачей.}%
  {\bfseries Текст ПЗ, презентация}%
  {\bfseries TODO: срок}{} }

\taskliterature{
\nocite{Devlin2019Bert}
\nocite{Reimers2019SentenceBert}
\nocite{Manning2008IR}
\nocite{JurafskyMartin2024}
\nocite{WeaviateDocs}
\nocite{MongoDBDocs}
\nocite{DjangoDocs}
\nocite{ISO25010}
\nocite{UMLSpec}
}

\authortaskapproval{}%
\supervisortaskapproval{}%

\authorrspzapproval{}%
\supervisorrspzapproval{}%
\consultantrspzapproval{}%

\supervisorrspzgrade{}%

\authorpzapproval{}%
\supervisorpzapproval{}%
\consultantpzapproval{}%

\supervisorpzgrade{}%

%%% Local Variables:
%%% TeX-engine: xetex
%%% eval: (setq-local TeX-master (concat "../" (seq-find (-cut string-match ".*-1-task\.tex$" <>) (directory-files ".."))))
%%% End:


\newcounter{totalfigures}
\newcounter{totaltables}
\newcounter{totallistings}

%%% Local Variables:
%%% TeX-engine: xetex
%%% eval: (setq-local TeX-master (concat "../" (seq-find (-cut string-match ".*-3-pz\.tex$" <>) (directory-files ".."))))
%%% End:
